\section{Discussion}
\label{sec:discussion}

\paragraph{Why can a change of about one intensity level flip an episode?} The open-loop check bears on three explanations. First, the closed loop amplifies small differences: we find no association between a changed first action and a flip, while almost every OpenVLA-7B trajectory ends at a different step, so flips appear to be decided by differences that accumulate over the rollout. Second, OpenVLA-7B chooses each action dimension greedily among \Bins{} bins: nearly tied tokens switch most often, but confident tokens switch as well and switched tokens jump by tens of bins, so a step moves OpenVLA-7B between distinct actions rather than nudging it. Third, OpenVLA-OFT regresses continuous actions: under the three mild steps it moves by \OlOftFirstJpegMean{} to \OlOftFirstBilMean{} bins on first frames and under one bin along the rollouts, and only the crop step, which moves it by about ten bins, shows a candidate loss in the closed loop. These observations fit the explanations without proving them.

\paragraph{What this means for practice.} Matching the image steps is cheap, and our bounds do not license skipping them. For the two checkpoints studied here, the larger risk to comparisons is not one step: any change to the evaluation reshuffles the fragile episodes, and other studies show that for some stacks even a different machine does~\citep{jiang2026bench,roboscope}. Score gaps within the band of Equation~\eqref{eq:band} should not be read as differences between models or pipelines unless the episodes are paired and the flips are reported. For policies that sample their actions the floor is higher still: the LIBERO-CTRL repository reports that \ExtCtrlSampledLo\% to \ExtCtrlSampledHi\% of rollouts of sampling policies change outcome from one rollout to the next~\citep{liberoctrlrepo}, so an A/A floor has to be measured on the policy being compared, not borrowed from a deterministic one.

We therefore recommend that evaluations (i) state their image pipeline and match the original one; (ii) compare pipelines or models on paired initial states, with the environment loop matched so that paired episodes start from the same scene, and report lost and gained episodes; (iii) measure an A/A floor on the same hardware and policy; and (iv) read a score gap within $\pm \Zval\sqrt{f}/n$ as no evidence of a difference.

\paragraph{Limitations.} The study covers one suite in simulation, one released checkpoint per model, one simulator version and one GPU type, so nothing is claimed about real robots, other suites or other policies. On LIBERO-Spatial the instruction is largely redundant with the image~\citep{xie2026language}, so the steps may matter less on suites where the instruction carries the goal, such as LIBERO-Goal. Closed-loop success is a coarse, noisy measure: a step triggers flips whose direction does not follow the step, so the study bounds the effect of a step rather than measuring it, and \NPairs{} initial states are not enough to resolve \ThresholdPts{} points at this rate of flips. OpenVLA-OFT is close to full success, which hides losses. The candidate effect of the crop on OpenVLA-OFT rests on one task; a rerun of its cell on the CPU of the other OpenVLA-OFT cells repeats its outcomes (Appendix~\ref{app:cpucontrol}). with the control rerun.
The image-level measurements of Appendix~\ref{app:eda} use frames rendered on a CPU; the GPU frames of the runs agree with them to within a few hundredths of a level. The open-loop check is exploratory: its summaries were chosen after the closed-loop results, it covers the first frame of every initial state but rollouts of only \OlSevStates{} (OpenVLA-7B) and \OlOftStates{} (OpenVLA-OFT) initial states per task, and it measures OpenVLA-OFT in the bins of OpenVLA-7B, averaged over a chunk.

\section{Conclusion}
\label{sec:conclusion}

We changed the image steps of the official OpenVLA evaluation one at a time and measured two released OpenVLA-family checkpoints on all \NPairs{} initial states of LIBERO-Spatial, with paired episodes, a rerun as the noise floor and a primary analysis fixed before any variant was run. No step has a net effect that \NPairs{} pairs can resolve: for the primary comparison, turning off the JPEG round trip, the loss is at most \SevJpegUB{} points with \AlphaPct\% confidence, which does not rule out the \ThresholdPts{} points that matter, and the closest secondary effect, the center crop on OpenVLA-OFT, is exploratory. Yet every step, including a JPEG change of about one intensity level, flips \SevFlipPctMin\% to \SevFlipPctMax\% of OpenVLA-7B episodes in both directions and changes almost every trajectory, so the cancellation that paired evaluations have reported for large perturbations already occurs for the evaluator's own steps on clean frames. The same paired design applies to other suites, to more checkpoints per model and to other evaluation choices, such as the simulator version.
